%!TEX root = ./../../springer-mono-format/book.tex
%!TEX root = ./../../thesis-trier.tex

Here, \cref{constr_sync:thm:PSPACE} shows that \constrsync{L} 
is decidable and solvable in polynomial space for regular $L \subseteq \Sigma^*$.
Then, \cref{constr_sync:thm:introductory_examples} 
presents various example problems and their complexities.

\begin{theorem}
	\label{constr_sync:thm:PSPACE}
	Let $L \subseteq \Sigma^*$ be regular.
	Then \constrsync{L} is decidable and in \PSPACE.
\end{theorem}
\begin{proof}
	An input semiautomaton $\semiautomat{A}{Q}{\Sigma}{\delta}$
	admits a synchronizing word in~$L$, i.e. $L \cap \Syn(\mathcal A) \ne \emptyset$,
	if and only if there exists $t \in Q$ such that
	\[  
	L \cap \bigcap_{q \in Q} L(\mathcal A_{q, \{t\}}) \ne \emptyset.
	\]
	Let $\mathcal B$ be a complete automaton 
	with $L = L(\mathcal B)$, which can be chosen in advance
	since~$L$ is fixed by the definition of the problem.
	
	An algorithm solving the problem
	can guess the state~$t$
	and then call a procedure
	to solve \DFAINT for $\mathcal B$
	and the automata 
	$\mathcal A_{q, \{t\}}$ for $q \in Q$
	as input (so that it is called with $|Q|+1$ automata as arguments).
	By \cref{prelim:thm:complexities},
	\DFAINT is in \PSPACE.
	So this algorithm can run in nondeterministic polynomial space.
	Since $\PSPACE = \NPSPACE$ 
	by \cref{prelim:eqn:cc_inclusions}
	we obtain the claim.
\end{proof}

The following introductory examples (one for each complexity class that will come up in this part)
are presented here to motivate the later results and to have some easy and concrete cases
with self-contained and independent proofs. Each example will be generalized in the coming sections.

\begin{theorem}
	\label{constr_sync:thm:introductory_examples}
	Let $a,b,c \in \Sigma$ be three distinct letters.
	\begin{enumerate}
		
		\item \label{constrsync:ex:item:a} $\constrsync{\{a\}}$ is in $\AC^0$.
		
		\item \label{constrsync:ex:item:bstar} $\constrsync{b^*}$ and $\constrsync{ab^*}$ are \LOGSPACE-complete
	    even for trap-state input automata.
		
		\item \label{constrsync:ex:item:ab} $\constrsync{(a+b)^*}$ and $\constrsync{\Sigma^*}$ are \NL-complete. 
	
		\item \label{constrsync:ex:item:abstara} $\constrsync{ab^*a}$ is \NP-complete even for input automata with a trap state~\cite{FernauGHHVW19}.
		
		\item \label{constrsync:ex:item:abcL} $\constrsync{(a+b)^*c^lL}$ is \PSPACE-complete for 
		non-empty and regular $L \subseteq \Sigma^*$ such that $|\fact{L} \cap \{a^n, b^m, c^l \}| \le 2$
		for some $n,m, l \ge 1$ even for trap-state automata as input.
		
		\item \label{constrsync:ex:item:cab} $\constrsync{c(a+b)^*}$ is \PSPACE-complete even for strongly connected automata as input~\cite{FernauGHHVW19}.	
	\end{enumerate}
\end{theorem}
\begin{proof}
	\ref{constrsync:ex:item:a}: 
	The letter $a$ synchronizes a given input automaton
	if and only if the transition matrix $A = (a_{ij})_{1 \le i,j \le n}$ 
	of $a$ contains a column of ones.
	Equivalently, there exists $J \in \{1,\ldots,n\}$
	such that for all $i \in \{1,\ldots,n\}$ we have 
	$a_{iJ} = 1$.
	This can be checked by an~$\AC^0$ circuit implementing the formula 
	\[
	\bigvee_{J \in \{1,\ldots,n\}} \left( \bigwedge_{i \in \{1,\ldots,n\}} a_{iJ} \right).
	\]
	In \cref{sec:constr_sync:finite} (\cref{constr_sync:thm:finite}) 
	this is generalized to finite constraint languages.
		
	\ref{constrsync:ex:item:bstar}:  
	In \cref{thm:SYNC-solv} we have shown that \DFAST for unary
	automata is \LOGSPACE-complete. This problem is equivalent
	to $\constrsync{b^*}$.			
	By adding self-loops to each state for the letter $a$
	in an input automaton $\semiautomat{A}{Q}{\{b\}}{\delta}$
	over~$\{b\}$, we obtain 
	an $\leACZEROm$-reduction from this problem. So, $\constrsync{b^*} \le_m^{\text{AC}^0} \constrsync{ab^*}$.	
	Next, we get an $\leACZEROm$-reduction
	from $\constrsync{ab^*}$
	to \SFS over the unary alphabet $\{b\}$ 
	in the following way: for $\semiautomat{A}{Q}{\{a,b\}}{\delta}$
	construct the \SFS-instance $(\mathcal A_{|\{b\}}, S)$
	with $S = \applytrans{\delta}{Q}{a}$
	where w.r.t.~$Q = \{q_1, \ldots, q_n\}$
	a characteristic bit-vector of $S$ can 
	be computed by an~$\AC^0$ circuit 
	as, for $[a] = (m_{ij}^{(a)})_{i,j\in \{1,\ldots,n\}}$,
	we have $q_j \in S$ if and only if $m_{ij}^{(a)} = 1$ for some $i \in \{1,\ldots,n\}$.	
	Then the claim follows as \SFS over $\{b\}$ is in \LOGSPACE
	by \cref{thm:SF-ST-comm}. 
		
	\ref{constrsync:ex:item:ab}: $\constrsync{(a+b)^*}$ and $\constrsync{\Sigma^*}$ are
	equivalent to the unconstrained synchronization problem
	over $\{a,b\}$ and~$\Sigma$, which are \NL-complete by Theorem~\ref{prelim:complexity:general_aut_sync}
	since $|\Sigma| \ge 2$.
	
	\ref{constrsync:ex:item:abstara}:
	We give a reduction from \DST
	over a unary alphabet, which is \NP-complete by Theorem~\ref{compl:thm:comm_ST}.
	Let $(\mathcal A, S, T)$ be an instance
	with $\semiautomat{A}{Q}{\{b\}}{\delta}$ and $S, T \subseteq Q$
	such that $S \cap T = \emptyset$.
	Pick some $\hat s \in S$
	and let $\overline{S} \colon ( Q \setminus T ) \to S$
	be given by $\overline{S}(q) = \hat{s}$ if $q \in Q \setminus (S \cup T)$
	and $\overline{S}(q) = q$ otherwise. Construct
	$\semiautomat{A'}{Q \cup \{t\}}{\{a,b\}}{\delta'}$ with a new state $t \notin Q$ 
	where, for $q \in Q$ and $x \in \{a,b\}$,
	\[
	\applytrans{\delta'}{q}{x} = \begin{cases}
		\applytrans{\delta}{q}{b} & \text{if } x = b, \\
		\overline{S}(q)           & \text{if } x = a \text{ and } q \in Q \setminus T, \\
		t                         & \text{if } x = a \text{ and } q \in T
	\end{cases}
	\]
	and $\applytrans{\delta'}{t}{x} = t$ for all $x \in \{a,b\}$. 
	The state $t$ must be the unique synchronizing state since it is a trap state.
	Furthermore, for $m \ge 0$,     
	we have $\applytrans{\delta}{S}{b^m} \subseteq T$
	if and only if $\applytrans{\delta'}{Q\cup \{t\}}{ab^ma} = \{t\}$,
	which yields that this is a valid reduction.
	
	In \cref{constr_sync:sec:motivation} we have already established that
	$\constrsync{ab^*a}$ is in \NP\ 
	by using the matrix orbit problem.
	For a different argument, given $\semiautomat{A}{Q}{\Sigma}{\delta}$,
	by \cref{prelim:word_in_PA} applied to $S = \delta(Q, a)$
	there exists a word in $ab^*a$ synchronizing~$\mathcal A$
	if and only if there exists $0 \le i < 2^{|Q|} - 1$
	such that $|\applytrans{\delta}{Q}{ab^ia}| = 1$.
	Using a binary encoding, $i$ can be guessed
	in polynomial time. By a square-and-multiply
	algorithm the transition matrix of $b^i$ can be computed in polynomial time. Then,
	the condition $|\applytrans{\delta}{Q}{ab^ia}| = 1$  
	can be checked in polynomial time.
		
	\ref{constrsync:ex:item:abcL}:
	By \cref{constr_sync:thm:PSPACE} we only have to establish \PSPACE-hardness.
	We give a reduction from \SIS over a binary alphabet, which
	is \PSPACE-complete by Theorem~\ref{prelim:thm:binary_pspace_problems}. 
	Let $(\mathcal A, S)$ be an instance of \SIS
	where $\semiautomat{A}{Q}{\{a_1, a_2\}}{\delta}$ 
	and $S \subseteq Q$. 		
			
	Let $u_1 = a^{n}b^{m}$, $u_2 = a^{n+1}b^{m}$.
	Intuitively, we re-encode the letters $a_1$ and $a_2$
	into words $u_1$ and~$u_2$, respectively, and build a new automaton~$\mathcal A'$
	containing a copy of~$Q$ in its state set 
	that branches along these two words in the same way as~$\mathcal A$
	branches on $a_1$ and~$a_2$.
	Additionally, the new automaton uses $c^l$
	to map states from~$S$ into a trap state.
	A word from the constraint language $(a+b)^*cL$
	synchronizes $\mathcal A'$
	if and only if, before reading $c^l$, we land
	inside~$S$. This follows since the suffix of this word following $c^l$ only uses letters
	from $\{u_1,u_2\}$, $\{u_1,c^l\}$, or $\{u_2,c^l\}$,
	but the construction is done in a way such that 
	we need all three words to make any changes to the current state set
	yielding a new synchronizing word.
	
	Set $P = \pref{\{a^n b^{m-1}, a^{n+1}b^{m-1}, c^{l-1} \}}$.
	Let $\semiautomat{A'}{ ( Q \times P ) \cup \{t\}}{\{a,b,c\}}{\delta'}$
	where, for $q \in Q$, $u \in P$,
	and $x \in \{a,b,c\}$, we set
	\[
	\delta'((q, u), x) =
	\begin{cases}
		(q, ux) & \text{if } ux \in P, \\
		(\delta(q, a_i), \varepsilon) & \text{if $ux = u_i$ for some $i \in \{1,2\}$,} \\	
		(q, \varepsilon) & \text{if $ux = c^l$ and $q \notin S$}, \\
		t                & \text{if $ux = c^l$ and $q \in S$},\\
		(q, \sigma_P(ux)) & \text{otherwise, i.e., if $ux \notin P \cup \{u_1, u_2, c^l\}$,}
	\end{cases}
	\]
	and $\delta'(t, x) = t$ for all $x \in \{a,b,c\}$.\footnote{Note that 
		for $u \in P$ and $x \in \{a,b,c\}$
		with $ux \notin P$
		we have $\sigma_P(ux) = x$ expect in the case $u = a^{n+1}$ and $x = a$.}

	\Cref{constrsync:fig:reduction} shows a part of the reduction
	for a single state~$q \in Q$ of the input automaton
	and $n = m = 2$.
	 The additional ``auxiliary'' states, used to track the current prefix and to branch when $abb$ or $aabb$ has been read, are indicated by rectangles. Transitions labeled~$c$ are omitted for clarity: they either return to $(q,\varepsilon)$ if $q \notin S$ or lead to the trap state~$t$. 
	 Along such a $c$-path $(q, c^i)$ with $1 \le i \le l-2$, 
	 reading $a$ or $b$ leads to $(q,\varepsilon)$.
	
		\begin{figure}[htb]
			\centering
			\begin{tikzpicture}				
				\node (q) {$(q, \varepsilon)$};
				\node[draw,rectangle, below = of q] (qa) {$(q, a)$};
				\node[draw,rectangle, below = of qa] (qaa) {$(q, aa)$};
				\node[draw,rectangle, below left = of qa] (qab) {$(q, ab)$};
				\node[below = of qab] (qabb) {$(\delta(q, a_1),\varepsilon)$};
				\node[draw, rectangle, below right = of qaa] (qaab) {$(q, aab)$};
				\node[below = of qaab] (qaabb) {$(\delta(q, a_2), \varepsilon)$};
				
				\node[right = of q, yshift=1em] (q-dummy) {$\ldots$};
				\node[right = of qa, yshift=1em] (qa-dummy) {$\ldots$};
				\node[right = of qaa, yshift=1.7em] (qaa-dummy) {$\ldots$};
				\node[right = of qaab, yshift=1em] (qaab-dummy) {$\ldots$};
				\node[left = of qab, yshift=1em] (qab-dummy) {$\ldots$};
				\node[left = of qabb, yshift=1em] (qabb-dummy) {$\ldots$};
				\node[right = of qaabb, yshift=1em] (qaabb-dummy) {$\ldots$};
				
				\path[->] (q) edge [right] node {$a$} (qa)
				(qa) edge [right] node {$a$} (qaa)
				(qaa) edge [below] node {$b$} (qaab);
				
				\path[->] (qa) edge [above] node {$b$} (qab)
				(qab) edge [left] node {$b$} (qabb);
				
				\path[->] (qaab) edge [right] node {$b$} (qaabb);
				
				\path[->] (qaab) edge [bend right=68,right] node {$a$} (q)
				(qab) edge [bend left,left] node {$a$} (q);
				\path[->] (qaa) edge [loop right] node {$a$} (qaa);
				\path[->] (q) edge [loop left] node {$b$} (q);
				
				\path[->] (q) edge [above] node {$c$} (q-dummy);
				\path[->] (qa) edge [above] node {$c$} (qa-dummy);
				\path[->] (qaa) edge [above] node {$c$} (qaa-dummy);
				\path[->] (qab) edge [above] node {$c$} (qab-dummy);
				\path[->] (qaab) edge [above] node {$c$} (qaab-dummy);
				\path[->] (qaabb) edge [above] node {$c$} (qaabb-dummy);
				\path[->] (qabb) edge [above] node {$c$} (qabb-dummy);				
			\end{tikzpicture}
			\caption[Construction for hardness of constraint language $\{a,b\}^*c^lL$]
			{Some states for an example of the construction in \cref{constrsync:ex:item:abcL} with $n = m = 2$,
				see the proof of \cref{constrsync:ex:item:abcL} for an explanation.}
			\label{constrsync:fig:reduction}
		\end{figure}
	
	Next, we show that $(\mathcal A, S)$ is a \rYES-instance of \SIS 
	if and only if $\Syn(\mathcal A')$ contains a word from the constraint language.
			
	First, suppose $\delta(Q, w) \subseteq S$ for some $w \in \{a_1,a_2\}^*$.
	We have 
	\[
	\delta'(( Q \times P ) \cup \{t\}, b^{m+1}) = Q \times \{\varepsilon\} \cup \{t\}.
	\]
	Let $\varphi \colon \{a_1,a_2\}^* \to \{a,b,c\}^*$
	be the morphism given by $\varphi(a_i) = u_i$ for $i \in \{1,2\}$.
	Since $\delta'((q, \varepsilon), u_i) = (\delta(q, a_i), \varepsilon)$
	for $i \in \{1,2\}$ and $q \in Q$,
	we find $\delta'((q, \varepsilon), \varphi(w)) = (\delta(q, w), \varepsilon)$
	for all $q \in Q$
	and thus $\delta'(Q \times \{\varepsilon\}, \varphi(w)) \subseteq S \times \{\varepsilon\}$.
	For $q \in S$
	we have $\delta'((q, \varepsilon), c^l) = t$.
	Since~$L$ is non-empty, we can pick some word $v \in L$.
	Combining everything and since~$t$ is a trap state,
	we find that 	
	$\delta'((Q \times P) \cup \{t\}, b^{m+1} \varphi(w) c^l v) = \{t\}$.
		
	Conversely, suppose we have $|\delta'(Q \times P \cup \{t\}, wc^lv)| = 1$
	for some $w \in \{a,b\}^*$ and $v \in L$.
	Since~$t$ is a trap state, we find $\delta'(Q \times P \cup \{t\}, wc^lv) = \{t\}$.
	Let $q \in Q$.
	We have $\delta'((q, \varepsilon), w) = (p, u)$
	for some $u \in \pref{\{u_1, u_2\}}$ and $p \in Q$. Assume $p \notin S$.
	Then $\delta'((p, u), c^l) = (p, \varepsilon)$ (this holds for $u = \varepsilon$ and 
	for $u \ne \varepsilon$).	
	Next, to map $(p, \varepsilon)$ to~$t$, we must first map it to a state
	in $S \times P$ and then apply $c^l$.	
	Now, for $w' \in \{a,b,c\}^*$
	to have $\delta'((p, \varepsilon), w') \in S \times P$
	it must have $u_1$ or $u_2$ as a factor,
	and to map a state from~$S \times P$
	to~$t$ it must have $c^l$ as a factor.
	Since by assumption $|\fact{L} \cap \{a^n, b^m, c^l\}| \le 2$,
	we have $c^l \notin \fact{L}$
	or $|\fact{L} \cap \{a^n, b^m\}| \le 1$.
	where the latter case yields $u_1, u_1 \notin \fact{L}$.
	Thus there is no way to map $(p, \varepsilon)$
	to~$t$ with a word from~$L$	and so we must have $p \in S$.
	Now, we have $\delta'((q, \varepsilon), w) = (p, u)$
	if and only if
	there exist words $v_0, v_1, \cdots, v_r$ 
	and $w_1, w_2, \cdots w_r \in \{u_1, u_2\}$
	such that
	\begin{enumerate}
	\item $w = v_0 w_1 v_1 w_2 \cdots w_r v_r$,
		
	\item $v_r$ has $u$ as a suffix, and
	
	\item $\delta'((q, \varepsilon), w_1 \cdots w_r) = (p, \varepsilon)$. 	
	\end{enumerate}
	Let $x_i \in \{a_1, a_2\}$ for $i \in\{1,\ldots,r\}$
	with $x_i = a_1$ if $w_i = u_1$
	and $x_i = a_2$ otherwise.
	Then $\delta(q, x_1 \cdots x_r) = p$
	and thus $\delta(Q, x_1 \cdots x_r) \subseteq S$.
	
	
	\ref{constrsync:ex:item:cab}: By \cref{constr_sync:thm:PSPACE} we only have to establish \PSPACE-hardness.
	We give a reduction from \SFS over a binary alphabet, which is \PSPACE-complete
	by \cref{prelim:thm:binary_pspace_problems} even for strongly connected automata.
	Let $(\mathcal A, S)$ be
	an \SFS-instance with $\semiautomat{A}{Q}{\{a,b\}}{\delta}$
	and $S \subseteq Q$. Pick some $\hat s \in S$.
	By extending~$\delta$ with $\applytrans{\delta}{q}{c} = q$ for each $q \in S$
	and $\applytrans{\delta}{q}{c} = \hat s$ for $q \notin S$
	we obtain an automaton over $\{a,b,c\}$ with 
	$\applytrans{\delta}{Q}{c} = S$
	that has a synchronizing word of the form~$cu$ for $u \in \{a,b\}^*$
	if and only if $|\applytrans{\delta}{S}{u}| = 1$.	
\end{proof}